This study asked whether early diffusion structure carries useful information beyond early volume and speed, and whether that information differs across two tasks: veracity screening and eventual reach prediction. The revised design tests this question with FibVID as the main dataset, Twitter15/16 as a comparative benchmark, size-based windows as the main specification, and time-window results as a robustness check.

\subsection{What the Results Show}
The clearest pattern in the revised results is still task asymmetry. On FibVID, veracity performance is above chance but modest throughout the main $k$-window analysis. The best baseline logit result reaches AUC $\approx 0.591$ at $k=180$, while richer bundles do not produce consistent gains over that baseline. In that sense, early diffusion signals do contain some information about misinformation risk, but they do not support a strong classification claim.

For reach, the pattern is much stronger. Under the harmonized FibVID specification, prediction performance is high even at moderate windows and remains strong in later ones, with OLS reaching about $R^2=0.886$ at $k=180$ and the interaction-full specification peaking around $R^2=0.861$ at $k=120$. This suggests that early diffusion structure is much more informative for \emph{how far} a cascade will spread than for \emph{whether} it is false.

The Twitter15/16 benchmark supports the same broad reading. In the original rumor-tree setting, veracity reaches about AUC $=0.657$ at $k=20$, while reach reaches about $R^2=0.719$ at $k=180$. Relative to that benchmark, FibVID weakens the veracity side but strengthens the reach side. The comparison therefore does not overturn the earlier asymmetry; it makes it clearer.

This point also matters for how the results speak back to the literature. Work such as \citet{vosoughi_spread_2018} shows that true and false information can diffuse differently when cascades are observed in full. The present results suggest a narrower conclusion under an early-warning constraint: visible diffusion differences in fully observed cascades do not translate automatically into strong early veracity classification. The reach result generalizes more clearly than the veracity result.

\subsection{Why the Methodological Changes Matter}
The residualization diagnostics (Results Table~\ref{tab:residualization_diag}) indicate that structural features were meaningfully volume-correlated before adjustment and nearly orthogonal to log early volume after adjustment. This matters because, without adjustment, apparent structure effects can be volume effects in disguise. Given how strongly early volume predicts later reach, that distinction is important.

This is also where the paper connects back to \citet{goel_structural_2016}. If structural quantities are mechanically tied to cascade size, then any predictive success can be hard to interpret. The residualization step does not solve that problem completely, but it does make the feature bundles more defensible by separating shape from early volume as much as the observed data allow.

Primary reliance on $W_k$ windows also remains important for interpretation. By fixing observed size, these analyses reduce confounding by raw arrival speed and make comparisons across cascades more structurally meaningful. The time-window results remain useful, but they answer a slightly different question. They show what can be recovered after a fixed amount of elapsed time, not after a fixed amount of observed information.

\subsection{Modeling Implications}
The model-family comparisons point to a similar task split. For veracity, logit remains competitive in the main FibVID analysis, and the gains from greater model flexibility are limited. This is consistent with the broader pattern in the paper: the veracity signal is present, but weak enough that more flexible models do not transform the task.

For reach, the situation is different. Tree-based regressors become very strong in larger windows, which suggests that non-linear interactions matter more for diffusion-outcome prediction than for misinformation-risk screening. In practical terms, this implies two different modeling regimes. For early veracity screening, relatively simple probabilistic models may already capture most of the available signal. For reach forecasting, richer interactions and nonlinearities are more important.

\subsection{Strength of Evidence and Inference Boundaries}
The permutation tests make the task asymmetry sharper. On FibVID at $k=60$, observed reach $R^2$ is far above the shuffled-label null, while veracity AUC is only modestly separated from it. The Twitter15/16 benchmark shows the same qualitative pattern, although the veracity signal is somewhat stronger there. These checks make it difficult to read the veracity results as a strong classification success.

At the same time, the time-window results do not reverse the main pattern. They are weaker than the size-based results, especially for reach, but they still point in the same direction: reach is more predictable than veracity, and $k$-windows remain the stronger main specification. I therefore interpret the time-window analysis as a robustness check on window definition rather than as a competing primary design.

\subsection{Limitations and Threats to Validity}
\paragraph{Representation dependence.}
The main FibVID analysis depends on a harmonized cascade-compatible representation. This is a transparent and rule-based conversion, but it is still a representation choice rather than a claim that FibVID is natively identical to Twitter15/16. The main results should therefore be interpreted as evidence under that representation. The supplementary graph-native veracity check is useful here: it shows that some conclusions, especially on the veracity side, are sensitive to how propagation structure is represented.

\paragraph{Predictive, not causal.}
Even after residualization, diffusion features can proxy unobserved event-level shocks, platform effects, or coordinated behavior. These estimates are predictive associations, not causal effects of structure on misinformation or reach.

\paragraph{Finite-fold uncertainty.}
Inference uses fixed cross-validation folds. Some late-window reach estimates are very strong, and while this is substantively informative, their exact magnitude should not be overinterpreted. The paper is more credible on the direction and relative strength of the effects than on any single point estimate.

\paragraph{Representation limits of interaction traces.}
Both the released rumor trees in Twitter15/16 and the propagation records in FibVID are interaction traces rather than full exposure networks. They do not observe recommendation systems, passive exposure, or off-platform diffusion. The models therefore recover signal from visible propagation, not from the full information environment in which these cascades spread.

\subsection{Implications for Early Warning Systems}
The practical implication is a staged early-warning logic. For veracity, diffusion-only performance remains limited in both datasets, so the most defensible use is early screening or triage rather than confident truth classification. Such signals may still help prioritize which cascades deserve further verification, but they should be interpreted conservatively.

For reach, the case is much stronger. As more early interactions accumulate, structural and dynamic features become highly informative about eventual spread. This makes diffusion-based models more naturally suited to prioritizing likely high-impact cascades than to deciding whether a claim is true or false.

Methodologically, the main implication is that diffusion-based early warning should not be framed as a single undifferentiated problem. The results suggest that task, window definition, and structural representation all matter. In this paper, the most defensible main design is a size-based, cascade-compatible specification for the reach task, with more cautious interpretation on the veracity side.

\subsection{Data and Code Availability}
Replication code for all analyses and figures is available at
\href{https://github.com/JiahaoZhang2001/Early-diffusion-signals.git}{\texttt{JiahaoZhang2001/Early-diffusion-signals}}.
This study uses publicly released rumor-cascade datasets, including FibVID and Twitter15/16.
To respect platform terms and redistribution constraints, I do not re-host raw tweet
content. The repository includes preprocessing and feature-extraction scripts,
along with environment requirements and run instructions to reproduce the results.
